\subsection{The precise target on a fixed cut stratum}

The next result gives the radius restriction for every stable type and
computes the two maps proposed on 67/081--082. We use the radius theorems
proved in Section~\ref{C:A2b-energy-section} and the subsequent image
theorem. Thus the restriction of the target is determined intrinsically,
not by an unspecified ``sufficiently small'' condition. Here a labelled
cut surface means the moduli-groupoid object with its cut circles and
punctures labelled, not a surface with an additional global topological
marking. In the fully marked Teichmuller space a boundary twist is real,
not circle-valued. One then lifts each tensor phase to its universal
cover and obtains the corresponding statements there. Thus the
integral Dehn twists are quotiented on both sides, or retained on both
sides; they are not silently lost from only one side.

\begin{theorem}[Proved: the exact propagator domain and the tangent map]
\label{A:graph-radius-domain}
Let $\Gamma$ be a connected stable graph of type $(g,n)$, with edge set
$E\ne\varnothing$. Fix its pointed compact normalization components
$X_v$, with every branch and external marking labelled. For an ordinary
fibre, retain an identification of the normalization with these fixed
objects; merely labelling the special points need not remove their
finite automorphisms. All constructions below are equivariant for
those automorphisms and also define the corresponding groupoid maps.
We first take
all external ends to be cusps. At each special point choose a local
coordinate, used only to measure radii and tensor moduli. Let $I_v$ be
the set of node branches on $X_v$. For $a\in[0,\infty)^E$, define
$(B_va)_h=a_e$ when the branch $h$ belongs to the edge $e$; both branches
of a loop are included separately. Let $H_v$ be the Jenkins--Strebel
support function for these active points, with all other special points
kept marked, and put
\[
 H_\Gamma(a)=\sum_v H_v(B_va),\qquad H_v(0)=0.
\]
Define the explicit open convex set
\[
 K_\Gamma=\left\{y\in\R^E:
       \sum_e a_e y_e<H_\Gamma(a)
       \text{ for every }a\geq0,\ \sum_ea_e=1\right\}.
\]
The attainable nonzero node tensors are exactly those with
$(\log|q_e|)_e\in K_\Gamma$. Every such tensor system has a unique
labelled geodesic-cut preimage, up to the equivalence of
Theorem~\ref{A:fullcaps}. More precisely, if
$x_e=(\ell_e/2\pi)^2$, the logarithmic modulus map is
\[
 G_\Gamma(x)=\sum_v B_v^T F_v(B_vx),\qquad
 G_\Gamma:(0,\infty)^E\xrightarrow{\ \sim\ }K_\Gamma,
\]
where $F_v$ is the vector of logarithmic cap radii on $X_v$.
It is a real analytic diffeomorphism with symmetric positive definite
derivative. In the phase coordinates specified by normalized cap maps,
the complete tangent map is block diagonal with blocks $DG_\Gamma$
and the identity on $\R^E$. This is the tangent map on the fixed
normalization fibre. When the normalization components also vary,
the differential is block triangular with the identity on the
normalization-moduli base and this invertible fibre block. In
particular the full real tangent map is invertible everywhere on
the fixed graph stratum.
\end{theorem}

\begin{proof}
For positive $x$, Mondello's inverse provides on each $X_v$ the unique
cap system with branch lengths $2\pi\sqrt{x_e}$. The two lengths paired
at an edge agree by construction. Lemma~\ref{A:tensor} gives exactly the
displayed formula for $G_\Gamma$, since the logarithm of a product of
the two radii is their sum. Conversely every geodesic-cut source has
these equal branch lengths, and hence is obtained by this construction.

The chain rule and the energy identity give
\[
 DG_\Gamma(x)=\sum_v B_v^T DF_v(B_vx)B_v.
\]
For any nonzero $u\in\R^E$, some branch of an edge with $u_e\ne0$
has nonzero coordinate in $B_vu$. The corresponding positive definite
summand makes $u^TDG_\Gamma(x)u>0$. Vertices without node branches
contribute zero and cause no exception. Integrating this inequality
along line segments proves injectivity, and the inverse function
theorem gives a local real analytic inverse.

The componentwise strict support inequalities imply
$G_\Gamma(x)\in K_\Gamma$: a nonzero $a$ has at least one nonzero
$B_va$, at which the inequality is strict. The finite-area support
gap on the components also gives the uniform bound
\[
 0<H_\Gamma(x)-x\cdot G_\Gamma(x)
       \leq\sum_v(2g_v-2+n_v)=2g-2+n.
\]
Here $n_v$ counts all node branches and external markings. The last
identity follows from $g=\sum_vg_v+|E|-|V|+1$ and
$\sum_v n_v=2|E|+n$.

We give the properness argument explicitly. Suppose $G_\Gamma(x_m)$
converges to $y\in K_\Gamma$. If $s_m=\sum_e x_{m,e}$ tends to
infinity along a subsequence, pass to another subsequence with
$a_m=x_m/s_m\to a$, $a\geq0$, $\sum a_e=1$. Homogeneity and the gap
give $H_\Gamma(a_m)-a_m\cdot G_\Gamma(x_m)\to0$.
Continuity of the support functions therefore gives
$H_\Gamma(a)=a\cdot y$, contradicting $y\in K_\Gamma$.
If a bounded subsequence has $x_{m,e}\to0$, the uniform collar bound
at either branch makes its log radius tend to $-\infty$. The other
branch radius has a fixed upper bound on its fixed normalization
component, so $(G_\Gamma(x_m))_e\to-\infty$, again a contradiction.
The inverse image of every compact subset of $K_\Gamma$ is therefore
compact in the positive orthant. Its image is nonempty, open, and
closed in the connected set $K_\Gamma$, and hence equals $K_\Gamma$.
This proves the asserted global inverse.

Choose the cap parametrizations so that their derivatives in the
specified local coordinates are positive real numbers. Each edge's
gluing phase $e^{i\theta_e}$ is then precisely the phase of $q_e$.
The source coordinates $(x,\theta)$ therefore map to
$(G_\Gamma(x),\theta)$. Their differential has the claimed blocks.
The welding equivalence supplies the surface and its uniqueness.
Changing coordinate norms translates the logarithmic target and its
support inequalities by the same vector, and changing coordinate
arguments rotates the phase torsors. Thus the intrinsic statement
does not depend on these auxiliary choices.
\end{proof}

\begin{corollary}[Proved: the fundamental-group test does not establish surjectivity]
\label{A:graph-pi-one}
For a fixed labelled normalization with its identification retained,
the map on ordinary fundamental groups
from the geodesic-cut source to the unrestricted tensor target is
the identity $\Z^E\to\Z^E$ in the twist generators, although its
image is the proper domain of Theorem~\ref{A:graph-radius-domain}.
The analogous statement holds for orbifold fundamental groups when
the normalization components vary and their finite automorphisms
are retained: inclusion of the restricted tensor domain is a
homotopy equivalence on each fixed graph stratum. This is compatible
with the everywhere invertible real tangent map. Neither test supplies
the missing properness onto the unrestricted target.
\end{corollary}

\begin{proof}
At a fixed normalization with its identification retained, take
$y_0=G_\Gamma(\mathbf 1)$. Linear
contraction in logarithmic radii to $y_0$ is a deformation retraction
of both $K_\Gamma\times(S^1)^E$ and
$\R^E\times(S^1)^E$ onto $\{y_0\}\times(S^1)^E$; convexity keeps
the former contraction inside $K_\Gamma$. Inclusion is therefore a
homotopy equivalence and preserves every circle generator. A full
turn in the normalized boundary twist is exactly a full turn in the
tensor phase, proving the assertion about the displayed generators.

Over varying normalization components use, intrinsically, the modulus
of the tensor produced when all node lengths equal $2\pi$ to define
the radial centre. This is a positive norm in each nodal tensor line;
it does not choose a phase in that line. The radial homotopies just
described are consequently defined without a local coordinate or
trivialization. Real analytic dependence in the component moduli
makes them continuous. They commute with relabellings and finite
automorphisms, and hence descend to the moduli orbifold, or equally
to its groupoid classifying space. The same deformation retraction
on both sides proves the orbifold assertion.

Properness fails for the larger target: take $x_m=t_m\mathbf1$ with
$t_m\to\infty$. The Jenkins--Strebel limit on every component shows
that $G_\Gamma(x_m)$ converges to a finite boundary point of
$K_\Gamma$, with equality in its positive diagonal support
inequality. Keeping all phases fixed gives a convergent sequence
of target tensors with no convergent source sequence. Thus the
inference from an invertible tangent map and a fundamental-group
isomorphism to a global isomorphism has a concrete failed hypothesis.
\end{proof}

\begin{example}[Proved: the entire key fibre in type $(0,4)$]
\label{A:04-exact-image}
Take the two labelled three-pointed spheres used in
Theorem~\ref{A:04}, with all four external ends cuspidal. Each node
cap has the same radius $r(\lambda)$, which increases strictly from
$0$ to $4$ as $\lambda$ increases from $0$ to infinity. The precise
tensor image in the specified coordinates is
\[
                    0<|q|<16.
\]
The parameter relation is $|q|=r(\lambda)^2$ and its phase is the
boundary twist. Thus this fibre is a punctured disc, whereas the
proposed unrestricted target fibre is $\C^*$. The derivative never
vanishes and both fundamental groups are $\Z$, with degree one map.
\end{example}

\begin{proof}
For one active puncture at zero on the three-pointed sphere,
the extremal Jenkins--Strebel domain is
$\C\setminus[1,\infty)$ and has conformal radius $4$, as the explicit
map $z=4u/(1+u)^2$ shows. The one-dimensional image theorem therefore
gives $F((0,\infty))=(-\infty,\log4)$, and the energy identity gives
$F'>0$. The two components are identically normalized, so their
radii agree at the common boundary length. Multiplication gives
exactly the claimed interval of tensor moduli. The assertions about
the derivative and fundamental group are also immediate from the
two preceding results.
\end{proof}

\begin{proposition}[Proved: the bordered image with all compatibility data retained]
\label{A:bordered-radius-algorithm}
The same finite-dimensional inversion determines the image when
external caps are retained. Add one positive variable $x_i$ for each
positive external boundary, and one output coordinate $\log r_i$
for its shadow. In $B_v$, copy a node variable to each of its two
flags and an external variable to its single flag. Then
$B^T(F_v)_v$ is a real analytic diffeomorphism onto the strict
support domain with support $\sum_vH_v(B_va)$.

For specified normalization components, node tensors, and actual
external pointed discs, the following is a necessary and sufficient
test for a geodesic-cut preimage: their tensor moduli and external
shadow radii belong to that domain; the unique recovered length
vector produces exactly the specified external discs. When both
conditions hold, the preimage is unique. Zero external radii are
handled on their corresponding face. This test states the additional
disc compatibility instead of assuming that forgetting node caps
preserves admissibility.
\end{proposition}

\begin{proof}
The matrix $B$ is still injective. The Hessian, support, finite-area
gap, and collar arguments in Theorem~\ref{A:graph-radius-domain}
therefore apply word for word, with a single cap radius instead of a
product at each external flag. This proves the indicated
diffeomorphism. A source necessarily gives its unique inverted
length vector and its canonical caps, proving necessity of both
tests. Conversely that length vector gives equal node-boundary
lengths, and the phase of each specified tensor gives its unique
toric gluing. The second test identifies the remaining external
caps with those required. Theorem~\ref{A:fullcaps} then supplies
the preimage. Its uniqueness follows from the radius inversion
and the uniqueness of toric gluing with its given phase.
\end{proof}
